% Cote 119: offprints, reports, courses, CV and letters (typed and printed part): archive cote 119, site page 211 (PDF page 212).
% typescript. Draft compared: none.
% AI-transcribed by Claude Opus 5.5 (claude-opus-5-5, adaptive thinking, effort high), Message Batches API, batch r119_full (msgbatch_01DwZNm5Y7YsKBWyAN4sFs8w), 2026-11-13; not yet checked by eye.
\scanpage{119/211}
\mspage{- 16 -}

M. HAKIM, Topos annelés et schémas relatifs, Ergebnisse des Math. Bd. 64, 1972,
Springer.

R. HARTSHORNE, Residues and Duality, Lecture Notes in Mathematics n° 20 (1966).

L. ILLUSIE, Complexe Cotangent et Déformations I, Lecture Notes in Math. n° 239
(1971), Springer et II, idem, n° 283 (1972).

R. KIEHL, Relativ analytische \unsure{Raüme|Raume}, Inventiones Math. vol. 16, pp. 40-112
(1972).

D. KNUTSON, Algebraic Spaces, Lecture Notes in Math. n° 203 (1971), Springer.

F.N. LAWVERE, Toposes, algebraic geometry and logic,Lecture notes in Math.,
n° 274 (1972), Springer.

I. MANIN, a) Théorie des groupes formels commutatifs sur des corps de caracté-%
ristique finie, (en russe) Uspekhi mat. Nauk, 1963, t. 18,
pp. 3-90 . (Il existe une traduction anglaise de  l'Amer. Math. Soc.).
b) Correspondances, motifs et transformations monoïdales (en russe),
Mat. Sbornik t. 77, pp. 475-507.

W. MESSING, The crystal associated to a Barsotti-Tate group, Lecture Notes in
Math. n° 264 (1971) Springer.

D. MUMFORD, Geometric Invariant Theory, Ergebnisse der Math. Bd 34, 1965,
Springer.

J.P. MURRE, On contravariant functors from the category of preschemes over a
field into the category of abelian groups, Pub. Math. n° 23,
pp. 5-43 (1964).

D. QUILLEN, Homotopical Algebra, Lecture Notes in Math. n° 43 (1967), Springer.

M. RAYNAUD, Spécialisation du Foncteur de Picard, Publications Math. n° 38,
pp. 27-76 (1970).

Mme M. RAYNAUD, Théorèmes de Lefschetz en cohomologie cohérente et cohomologie
étale, Thèse Paris 1972 (paraîtra dans Lecture Notes of Math.).

J.P. RAMIS et G. RUGET, Complexe dualisant et théorèmes de dualité en géométrie
analytique complexe, Pub. Math. n° 38, pp. 77 à 91 (1970).

\margin{211}
